% !TeX encoding = UTF-8
% Basic regression-style exercise of the simplified pxLST public API.
\documentclass{P3CTeX}

\usepackage{pxLST}

\begin{document}

\section*{pxLST basic happy-path usage}

% Global configuration: Darcula palette + Java listings, pxlst base style.
\pxCodeSet[dark]{pxJava}[pxlst]

Here we check that the main listing environments compile and render with
their expected signatures:

\begin{pxCode}[language=pxJava]{HelloWorld.java}
public class HelloWorld {
  public static void main(String[] args) {
    System.out.println("Hola, món!");
  }
}
\end{pxCode}

Inline code with \pxCodeIn[pxJava]{print("hola")} should compile inside
running text without breaking paragraphs.

% Framed variant with slightly different box styling.
\begin{pxCode*}[language=pxJava]{FramedHello.java}
class FramedHello {
  public static void main(String[] args) {
    System.out.println("Hola des de pxCode*");
  }
}
\end{pxCode*}

% CodeBox: side-note comment panel (non-breakable).
\begin{CodeBox}[language=pxJava]{bx:basic}{Hello world amb comentari}{nota}{Llistat bàsic amb nota al marge}
public class Boxed {
  public static void main(String[] args) {
    // Comentari bàsic
  }
}
\end{CodeBox}

% CodeBox* should be identical but breakable across pages; we do not force
% an actual page break here, only exercise the environment.
\begin{CodeBox*}[language=pxJava]{bx:break}{Hello world amb salt}{nota}{Versió breakable del mateix llistat}
public class BoxedBreak {
  public static void main(String[] args) {
    // Comentari llarg que podria forçar salts de línia
  }
}
\end{CodeBox*}

% File-based listing to exercise \pxCodeFile.
% The path is deliberately simple; users can adapt it to real project files.
\pxCodeFile[pxJava]{example-code/HelloWorld.java}

\end{document}

